\documentclass[english,twoside, openright, hidelinks, 11 pt, class=report,crop=false]{standalone}
\input{../../preamb}
\input{../bm}

\begin{document}
\opgt

\op{odpar} \vs
\abc{
\item Lag en rekursiv og en eksplisitt formel positivt partall nummerr $i$.
\item Lag en rekursiv og en eksplisitt formel for positivt oddetall nr $i$. 
}

\op{opgfolari}
Gitt følgen
\[ 5, 15, 25,\,... \]
\abc{
\item Finn et rekursivt uttrykk for $a_i$.
\item Finn et eksplisitt uttrykk for $a_i$.
\item Hva er verdien til $i$ når $a_i=125$?
}

\op{eksar}
Finn det eksplisitte uttrykket til  den aritmetiske følgen når du vet at\os
\abc{
\item $ a_1=3 $ og $ a_4 = 30 $ \os 
\item $ a_1 = 5 $ og $ a_{11} = -25 $ \os
\item $ a_3 =14 $ og $ a_5=26 $ 
}

\op{eksgeo}
Finn det eksplisitte uttrykket til den geometriske følgen når du vet at\os
\abc{
\item $ a_1=\frac{1}{2} $ og $ a_2 = \frac{1}{6} $ \os 
\item $ a_1 = 5 $ og $ a_4 = 40 $
}

\op{opgfolgeoi}
Gitt følgen
\[ 2, 8, 32,\,... \]
\abc{
	\item Finn et rekursivt uttrykk for $a_i$.
	\item Finn et eksplisitt uttrykk for $a_i$.
	\item Hva er verdien til $i$ når $a_i=128$?
}

\begin{comment}
	\op
Bruk formelen fra oppgave \textsl{\ref{sumkvad}a} og det eksplisitte uttrykket for en aritmetisk følge til å forklare at summen av en aritmetisk rekke kan skrives som:
\[na_1 +\frac{dn(n-1)}{2} \]


\op
a) Bruk (\ref{sumg}) og vis at summen $ S_\infty $ når $ n\to\infty $ blir:
\[ S_\infty=a_1\frac{1}{1-k} \]
for $ |k|<1 $.
\end{comment}
\nes

\op{parodd} \vs
\abc{
\item Bruk figuren under til å forklare at summen $ S_n $ av de $ n $ første naturlige tallene er gitt ved
\[S_n=\frac{n(n+1)}{2}  \]
\fig{sum}
\item Skriv opp summen av det første, de to første og de tre første  oddetallene. Bruk en lignende figur som i oppgave a) til å vise at summen $ S_n $ av de $ n $ første oddetallene er
\[ S_n = n^2 \] 
}
\op{sum10ar}
Finn $ S_{10} $ for rekkene:\os

\abc{
\item $ 7+13+19+25+\ldots $\quad
\item $ 1+9+17+25+\ldots $
}

\op{ar435}
Gitt rekken 
\[ 8+11+14+\ldots \]
Finn $ n $ er summen av rekken lik 435?

\op{opgarekfraeks}
Gitt rekken
\[ 3+7+11+... \] 
Finn $ n $ når $S_n=903$.

\op{opggerekfraeks}
Gitt rekken
\[ 3+6+12+24+... \]	
Finn $n$ når $S_n=93$?

\eksop{R2H25}{r2h25d1opg6a}
Finn summen til rekken
\[ -3+0+3+...+69 \]

\op{geon}
Gitt rekken
\[ 3+12+48+\ldots+768 \]
Finn summen av rekken. 
\newpage
\op{geoa12}
En geometrisk rekke har $ {a_1 = 2} $ og $ {k=3} $.\os 
\abc{
\item Vis at summen $ S_n $ kan skrives som:
\[ S_n = 3^n-1 \]
\item Regn ut summen for de tre første leddene.

\item For hvilken $ n $ er $ S_n=728 $?
}


\begin{comment}
\textbf{c)} Hvis du fortsetter å spare slik, og medregner innskudd samme måned, når vil du ha 24200 kr på konto? 
\end{comment}

\op{1over4}
Gitt den uendelige rekken
\[ 4+1+\frac{1}{4}+\ldots \]
\abc{
\item Forklar hvorfor rekken er konvergent.
\item Finn summen av den uendelige rekken.
}

\eksop{R2H25}{r2h25d1opg6b}
Gitt den uendelige geometriske rekka
\[ 5+5\left(\frac{1}{2}-x\right)+5\left(\frac{1}{2}-x\right)^2+... \]
Finn konvergensområdet til rekka.


\op{opggerekuendeks}
Gitt den uendelige rekken 
\[ 1+\frac{1}{x}+\frac{1}{x^2}+.... \]
\abc{
\item For hvilke $ x $ er rekken konvergent?
\item Vis at \y{S_\infty=\frac{x}{x-1}} når rekken konvergerer.
\item Finn $ x $ når $ S_\infty=\frac{3}{2} $.
\item Finn $ x $ når $S_\infty= -1 $.
}

\begin{comment}
	\op{stav} 
	Tenk at uendelig mange personer skal sette sammen en stav. Første person legger på en meter, neste person legger på 0.1 m, neste legger på 0.01 m osv. Hvor lang blir staven?
\end{comment}
\newpage
\op{099er1}
\abc{
\item Skriv desimaltallet 0.999... som en uendelig geometrisk rekke.
\item Forklar hvorfor rekken er konvergent og bruk dette til å finne summen av rekken. 
}

\op{geokonv}
Gitt den uendelige rekken
\[\frac{1}{3} +\frac{1}{3}(x-2)+ \frac{1}{3}(x-2)^2+\ldots\]
\abc{
\item For hvilke $ x $ er rekken konvergent? 

\item For hvilken $ x $ er $ S_n = \frac{2}{9} $?

\item For hvilken $ x $ er $ S_n= \frac{1}{6}$?
}

\nes
\op{ind}
Vis ved induksjon at for alle $ n\in \mathbb{N} $ er\os
\abc{
	\item $ 1+2+3+\ldots+n = \dfrac{n(n+1)}{2} $
	\item $ 1+2 +2^2 +\ldots+ 2^{n-1}= 2^n-1$
	\item $ 4+4^2+4^3+\ldots+4^n = \dfrac{4}{3}(4^n-1) $
	\item $ 1^2 + 2^2+3^2+ \ldots+ n^2 = \dfrac{n(2n+1)(n+1)}{6} $ 
}

\op{div3}
Vis ved induksjon at $ n(n^2+2) $ er delelig med 3 for alle $ n\in\mathbb{N} $.

\op{factorials}
\abc{
\item
Vis ved induksjon at:
\[\frac{1\cdot2}{1}\cdot\frac{1\cdot2\cdot3\cdot4}{1\cdot2\cdot3}\cdot
\ldots\cdot\frac{(2n)!}{(2n-1)!}=2^n n! \]
\textsl{Hint}: \y{(2(k+1))!=(2k+1)!(2k+2)}.\os

\item Hvordan kan venstresiden i a) skrives enklere? Utfør induksjonsbeviset på nytt etter forenklingen.
}
\newpage
\gruboprecc{opgfoliplusone}
Omskriv de rekursive og eksplisitte formlene fra \rref{arfol} og \rref{geofol} til formler som gir uttrykk for $a_{i+1}$ i steden for $a_i$. \vsk

\gruboprecc{opgfolarrekomv}
Forklar hvorfor den eksplisitte formelen til en aritmetisk rekke med $n$ element også kan skrives som
\[ a_i=a_n-(n-i)d \]

\grubeksoprecc{r2v26d1opg6}{R2V26}
\abc{
\item I en aritmetisk følge $\{a_i\}$ er $a_1=5$ og $d=-0,1$. Finn $i$ slik at følgen inneholder så mange elementer som mulig uten at noen elementer er mindre enn eller lik 0.
\item En figur består av uendelig mange kvadrater. La $l_i$ betegne sidelengden til kvadrat nummer $i$, hvor kvadratene er nummerert etter størrelse fra størst til minst. Da er $\frac{l_{i-1}-l_i}{l_{i-1}}=0,1$. Finn arealet til figuren når det største kvadratet har areal 19.
}

\grubeksoprecc{r2h24d1opg2}{R2H24} \vs
\abc{
\item Finn summen til den aritmetiske rekka $3+7+11+15+...+399$.
\item Bestem kvotienten $k$ til en konvergent uendelig geometrisk rekke som har $a_1=12$ og sum lik 18.
\item Vis at tallet $0,75757575...$ kan skrives som en uendelig\\ geometrisk rekke. Bruk dette til å vise at $0,75757575...=\frac{58}{33}$.
}

\grubeksoprecc{r2h24d2opg1b}{R2H24}
Gitt den uendelige geometrisk rekken $1+(\ln x - 1)+(\ln x - 1)^2+...$.\\
Avgjør om påstanden er rett. Grunngi svaret ditt.\\
\begin{center}
	\textit{Dersom} $x=\frac{1}{e}$\textit{, vil summen av rekka være} $\frac{1}{3}$.
\end{center}

\gruboprecc{eksmpr2v23d1opg4}
En aritmetisk rekke $a_1+a_2+...+a_n$ har sum $S_n=3n^2$ for $n\in\mathbb{N}$. \os Bestem $a_4$.

\gruboprecc{eksmpr2v23d1opg5}
En uendelig geometrisk rekke $b_1+b_2+...$ konvergerer mot 6. Summen av de tre første leddene er $\frac{38}{9}$.\os
 Bestem summen av de fire første leddene.
 
\grubeksoprecc{r2h23d1opg3}{R2H23D1}
En uendelig geometrisk rekke $ a_1+a_2+a_3+... $ konvergerer mot 8.
\abc{
	\item Bestem summen av de fire første leddene når du får vite at $ a_1=4 $. 
}
I en aritmetisk rekke er $ a_1+a_4+a_7 =114 $.
\abcs{2}{
	\item Bestem $ a_4 $. 
}

\gruboprecc{opgfolform} \vs
\abc{
\item Gitt $n,m\in\mathbb{N}$. Vis at vi for en aritmetisk følge $\{a_i\}$ har at 
\[ a_n=a_m+(m-n)\cdot d \]
hvor $d$ er den konstante differansen mellom to naboelement.
\item Gitt $n,m\in\mathbb{N}$. Vis at vi for en geometrisk følge $\{g_n\}$ har at 
\[ g_n=g_m k^{n-m} \]
hvor $k$ er den konstante kvotienten mellom to naboelement.
\item Løs oppgave \refop{eksgeo}c og \refop{eksar}b ved å bruke formlene fra a) og b).
}

\grubeksoprecc{r2v25d1opg3}{R2V25}
\abc{
\item Finn en rekursiv formel for følgen $2, 5, 9, 14, 20, ...$.
\item Vis ved induksjon at den eksplisitte formelen til følgen er gitt ved $a_i=\frac{i(i+3)}{3}$.
}


\gruboprecc{opgfol}
Desimaltallet 0.097097... er et desimaltall med repeterende desimalmønster. Desimalmønsteret 097 har verdi 97 og 3 repeterende siffer.
\abc{
\item Skriv 0,097097... som en geometrisk rekke
\item Skriv 0,097097... som en brøk.
\item La $a$ være et desimaltall som
\begin{itemize}
	\item er mindre enn 1
	\item har repeterende desimalmønster med verdi $b$ og $c$ repeterende siffer.
\end{itemize}
Vis at
\[ a=\frac{b}{10^c-1} \]
}
\mers{En nesten identisk oppgave er gitt i \mb, men der er den tenkt løst ved andre metoder enn her.}


\newpage
\grubop{opgfolg1o4expnsums}
Alle firkanter i figuren under er kvadrater. Bruk figuren til å argumentere for at
\[ \frac{1}{4}+\frac{1}{16}+\frac{1}{64}+...=\frac{1}{3} \]
\fig{opgfolg1o4expnsums}

\newpage
\grubop{opgfolgarchi}
I figuren under ligger alle markerte punkt på grafen til andregradsfunksjonen $f=ax^2+bx+c$. Tangenten til $f$ i $C$ er parallell med $AB$. Den samme relasjonen gjelder for $F$ og $AD$, $D$ og $AC$, $G$ og $DC$, $H$ og $CE$, $E$ og $CB$, og $I$ og $EB$. I en grubleoppgave i \tmen\, viste vi dette:
\st{
La $s$ være horisontalavstanden mellom $A$ og $B$.
\begin{itemize}
	\item Horisontalavstanden mellom $A$ og $C$ (og $C$ og $B$) er da $\frac{1}{2}s$.
	\item Arealet $T$ til $\triangle ABC$ er da $\frac{a}{8}s^3$.
\end{itemize}
}
\abc{
\item La $I$ være arealet avgrenset av grafen til $f$ og linjestykket $AB$. Videre tar vi det for gitt at $I$ er lik summen av de uendelig mange trekantene som kan dannes av linjestykker mellom tangeringspunkt (slik som den grønne trekanten, og de blå og de røde trekantene.). Vis at 
\[ I=\frac{1}{4}T+\frac{1}{16}T+\frac{1}{64}T+... \]
\item Vis at $I$ kan skrives som en uendelig geometrisk rekke. \\
\item Finn $I$ som en endelig sum uttrykt ved $T$.
}
\fig{opgfolgarchi}


\grubop{viseks3}
Bruk summen av en aritmetisk rekke til å vise at ligningen gitt i \textsl{Eksempel 3} på side \pageref{prodind} er sann.

\grubop{opgfolinduction}
Vis ved induksjon at for $n\in\mathbb{N}$ er
\[ \left(x^n\right)'=nx^{n-1} \]

\grubop{opgfolgeoarrek}
Gitt en følge med rekursiv formel
\[ a_i = ka_{i-1}+d \]
hvor $ k $ og $ d $ er konstanter. Finn en eksplisitt formel for følgen uttrykt ved $ a_1 $, $ k $, $ d $ og $ n\in\mathbb{N} $.

\grubop{opgsumnkvad}
Målet med denne oppgaven er å, uten bruk av induksjon, vise at summen av $ n $ kvadrater er gitt ved følgende formel:
\[ \sum\limits_{i=1}^n i^2 = \frac{n(2n+1)(n+1)}{6} \tag{I}\label{sumkvad}\]
\abc{
\item Forklar hvorfor vi kan skrive
\[ 1^2 + 2^2 + 3^2+\ldots = 1 + (1+3) + (1+3+5)+ \ldots  \]
\item Ut ifra det du fant i a), forklar at
\[ \sum\limits_{i=1}^n i^2 = n+\sum\limits_{i=1}^n (n-i)(2i+1)  \]
\item Skriv ut alle kjente summer fra b) og løs ligningen med\\ hensyn på $ \sum\limits_{i=1}^n i^2 $. Du skal da komme fram til (\ref{sumkvad}).
}

\grubop{opgfoldelelig}
Vis at hvis $b\in\mathbb{Z}/\{1\}$ og ${n\in\mathbb{N}}$, så er $b^n-1$ delelig med $b-1$.

\end{document}

